% =============================================================
% TMLR Cover Letter — PAPER6
% Capture-Resilient Voting in Multi-LLM Verifier Ensembles
% =============================================================
\documentclass[11pt]{letter}
\usepackage[margin=1in]{geometry}
\usepackage{hyperref}

\signature{Anonymous (for double-blind review)}
\address{Anonymous}

\begin{document}

\begin{letter}{TMLR Action Editor}

\textbf{Subject:} Submission --- ``An Empirical Study of Aggregation
Methods Under Verifier Capture in Multi-LLM Ensembles: When Does
Majority Voting Fail, and Can We Beat It?''

\bigskip

Dear Action Editor,

\bigskip

We are pleased to submit our manuscript for consideration by TMLR.
The paper presents an empirical benchmark of seven aggregation methods
for multi-LLM verifier ensembles under three verifier-capture attacks,
across four scenarios and three ensemble sizes.

\bigskip

\noindent\textbf{Central finding.} The headline result is
\emph{negative}: no aggregation method we tested substantially
outperforms plain majority voting under verifier capture, and several
trust-score variants \emph{underperform} majority under strong capture.
We argue this negative result is structural: a captured verifier is
statistically indistinguishable from an honest verifier that happens
to be wrong, so any rule that uses statistical deviation to identify
capture cannot outperform a rule that ignores the distinction.

\bigskip

\noindent\textbf{Why TMLR.} TMLR is the right venue because (i) the
contribution is a rigorous empirical benchmark with reproducible code,
not a conference-length systems paper; (ii) the negative result is
methodologically important for the growing community building LLM-judge
ensembles, and TMLR's scope explicitly covers empirical studies of ML
systems; and (iii) we release all code, attack implementations, and
analysis scripts under a permissive licence, consistent with TMLR's
reproducibility expectations.

\bigskip

\noindent\textbf{Contributions.}
\begin{enumerate}
\item A formal threat model for verifier capture, including a
      definition of $\varepsilon$-capture-resilience that decouples
      the adversary's per-verifier accuracy from its aggregate
      influence, plus a deterministic lower bound (Theorem 1) showing
      no rule can be $\varepsilon$-CR for $\varepsilon < 1/2$ when
      $f \geq N/2$.
\item A unified experimental protocol benchmarking seven aggregation
      methods (majority, PBFT threshold, uniform-weighted mean, three
      trust-score variants, and an adaptive hybrid) under three attack
      families (prompt injection, rubric injection, majority-bias
      mimicry) across $N \in \{5,7,10\}$ and $f \in \{1,2,3\}$, with
      4,000 trials per cell and 5 seeds.
\item A falsifiable prediction---that captured verifiers are
      statistically indistinguishable from honest-but-wrong verifiers---
      tested and supported across all scenarios, with nuanced exceptions
      (trust-score weighting beats majority under prompt-injection
      capture at low honest accuracy, but fails under majority-bias
      mimicry).
\item A small real-LLM validation pilot (8 LLM judges, 1,248 rubric
      decisions from a published benchmark) consistent with the
      synthetic-data trend; full real-trace benchmark is explicitly
      scoped as future work.
\end{enumerate}

\bigskip

\noindent\textbf{Reproducibility.} All code (crv.py, attacks.py,
extended\_methods.py, benchmark.py, verify\_p6.py), synthetic
generators, attack implementations, and analysis scripts are released
under a permissive licence with an attached ResponsibleUse.md file.
A full benchmark run (4 scenarios $\times$ 7 methods $\times$ 5 seeds,
$T=1000$ rounds) takes approximately 4 minutes on a single CPU core
and uses 50 MB of memory. SHA-256 hashes of all released artifacts
are recorded in HASHES.txt.

\bigskip

\noindent\textbf{Dual-use and responsible release.} The benchmark
implements three capture attacks; we have mitigated dual-use risk by
(i) releasing only defensive benchmark code, not captured-verifier
weights or fine-tuning scripts; (ii) attaching a ResponsibleUse.md
file prohibiting offensive use; and (iii) demonstrating in the paper
that these attacks have only small absolute effect against
well-calibrated honest majorities, limiting practical upside for
adversaries. The Broader Impact section discusses this in detail.

\bigskip

\noindent\textbf{Diff with prior submissions.} This is a fresh
submission. We have not submitted this work elsewhere.

\bigskip

\noindent\textbf{Recommended reviewers.} We suggest the AE draw from
the TMLR pool of reviewers with expertise in one or more of: LLM
judging / verifier ensembles, Byzantine fault tolerance, adversarial
robustness of ML systems, or calibration of probabilistic predictors.
We provide three suggested reviewers in a separate file (anonymized
for double-blind review).

\bigskip

\noindent\textbf{Conflicts of interest.} The authors declare no
financial conflicts of interest. The authors maintain an open-source
peer-review audit tool that is referenced as one of the motivating use
cases; the tool is cited consistently with the rest of the literature
and no special pleading is made for it.

\bigskip

We thank you for your time and look forward to your response.

\bigskip

Sincerely,

\bigskip

Anonymous (for double-blind review)

\end{letter}
\end{document}
